\documentclass[conference]{IEEEtran}

% -------------------------------------------------------
% Packages
% -------------------------------------------------------
\usepackage[T1]{fontenc}
\usepackage[utf8]{inputenc}
\usepackage{amsmath, amssymb}
\usepackage{graphicx}
\usepackage{booktabs}
\usepackage{hyperref}
\usepackage{cite}
\usepackage{newtxtext, newtxmath}
\usepackage[none]{hyphenat}

% -------------------------------------------------------
% Title
% -------------------------------------------------------
\title{Monte Carlo Simulation of Noise-Resilient Two-Qubit Variational Quantum Classifiers with Zero-Noise Extrapolation for Intent Classification}

% -------------------------------------------------------
% Author block
% -------------------------------------------------------
\author{
  \IEEEauthorblockN{John-Ronan S. Beira\textsuperscript{1 a *}}
  \IEEEauthorblockA{
    \textsuperscript{1}Bongao Alliance Church, Datu Halun Street, Bongao, Tawi-Tawi, Philippines\\
    \href{mailto:john-ronan.beira@g.msuiit.edu.ph}{\textsuperscript{a} john-ronan.beira@g.msuiit.edu.ph} / 
    \href{mailto:johnronanbeira100404@gmail.com}{johnronanbeira100404@gmail.com}
  }
}

% -------------------------------------------------------
% Abstract
% -------------------------------------------------------
\def\myabstract{%
  \mdseries
This study presents a Monte Carlo simulation framework to evaluate the
performance of a \textbf{2-qubit Variational Quantum Classifier (VQC)} for binary
intent classification under Noisy Intermediate-Scale Quantum (NISQ)
conditions. Using a focused subset of the 
\href{https://huggingface.co/datasets/tanaos/synthetic-intent-classifier-dataset-v1}{\textbf{Tanaos synthetic intent classifier v1}} dataset, 
the task is restricted to \textbf{``Greeting''} and \textbf{``Farewell''} categories to
suit the 2-qubit architecture. Text utterances are vectorized via TF-IDF and
reduced to a 2-dimensional feature space using Principal Component Analysis (PCA). 
The model architecture utilizes a ZZFeatureMap for quantum state
encoding and a RealAmplitudes variational circuit with six trainable
parameters. To combat hardware imperfections, we integrate \textbf{Zero Noise Extrapolation (ZNE)} 
as an error mitigation strategy. We systematically evaluate the pipeline by varying 
depolarizing noise levels (0--10\%), measurement shot counts (100--10{,}000), and circuit 
repetitions. Through 1,000 independent Monte Carlo trials per configuration, we report mean accuracy 
and statistical variance, showing that ZNE establishes a noise-resilient reliability baseline for 
near-term quantum NLU tasks.
}

% -------------------------------------------------------
% Keywords
% -------------------------------------------------------
\def\mykeywords{%
  \mdseries
  Variational Quantum Classifier (VQC), Monte Carlo Simulation,
  Noisy Intermediate-Scale Quantum (NISQ), Intent Classification (IC),
  Machine Learning, Quantum Machine Learning (QML), Quantum Noise,
  Tanaos Dataset%
}

% -------------------------------------------------------
% Bibliography
% -------------------------------------------------------
\def\bibPhukan{%
  A. Phukan and A. Ekbal,
  ``A survey of Quantum Natural Language Processing: From Compositional models to NISQ-Era Empiricism,''
  in \textit{Proc. QuantumNLP: Integrating Quantum Computing With Natural Language Processing},
  pp. 65--75, Mumbai, India, 2025.
  Association for Computational Linguistics.
  \url{https://doi.org/10.18653/v1/2025.quantumnlp-1.9}%
}

\def\bibZhu{%
  L. Zhu, Y. Dong, Z. Zhang, and X. Li,
  ``Rethinking Quantum Noise in Quantum Machine Learning: When Noise Improves Learning,''
  \textit{arXiv (Cornell University)}, 2026.
  \url{https://doi.org/10.48550/arxiv.2601.13275}%
}

\def\bibLung{%
  R. I. Lung and F. S. Duma,
  ``A Noisy Optimization Mechanism for Variational Quantum Classifiers,''
  \textit{Knowledge-Based Systems}, vol. 340, p. 115659, 2026.
  \url{https://doi.org/10.1016/j.knosys.2026.115659}%
}

\def\bibXhelollari{%
  F. Xhelollari, T. Li, and J. Chen,
  ``Robust Hybrid Quantum-Classical Latent Diffusion Models via Quantum Noise,''
  \textit{OpenReview}.
  \url{https://openreview.net/forum?id=NqTeVoz36v}%
}

\def\bibCincio{%
  L. Cincio, K. Rudinger, M. Sarovar, and P. J. Coles,
  ``Machine Learning of Noise-Resilient Quantum Circuits,''
  \textit{PRX Quantum}, vol. 2, no. 1, 2021.
  \url{https://doi.org/10.1103/prxquantum.2.010324}%
}

\def\bibHiggs{%
  ``Impact of Circuit Depth versus Qubit Count on Variational Quantum Classifiers for Higgs Boson Signal Detection,''
  \textit{arXiv}, n.d.
  \url{https://arxiv.org/html/2601.11937v1}%
}

\def\bibSeilkhan{%
  M. Seilkhan and A. Taizhanov,
  ``Comparing Classical and Quantum Variational Classifiers on the XOR Problem,''
  \textit{Open MIND}, 2026.
  \url{https://doi.org/10.48550/arxiv.2602.24220}%
}

\def\bibZengCoecke{%
  W. Zeng and B. Coecke,
  ``Quantum Algorithms for Compositional Natural Language Processing,''
  \textit{Electronic Proceedings in Theoretical Computer Science}, vol. 221, pp. 67--75, 2016.
  \url{https://doi.org/10.4204/eptcs.221.8}%
}

\def\bibPal{%
  S. Pal, P. Pakray, P. Jain, A. Ekbal, and S. Bandyopadhyay,
  ``Proceedings of the QuantumNLP: Integrating Quantum Computing with Natural Language Processing,''
  \textit{Proceedings of the QuantumNLP: Integrating Quantum Computing With Natural Language Processing}, 2025.
  \url{https://doi.org/10.18653/v1/2025.quantumnlp-1.0}%
}

\def\bibCoecke{%
  B. Coecke, D. F. Giovanni, K. Meichanetzidis, and A. Toumi,
  ``Foundations for Near-Term Quantum Natural Language Processing,''
  \textit{arXiv (Cornell University)}, 2020.
  \url{https://doi.org/10.48550/arxiv.2012.03755}%
}

\def\bibKrawchuk{%
  C. Krawchuk, N. Khatri, N. J. Ortega, and D. Kartsaklis,
  ``Efficient Generation of Parameterised Quantum Circuits from Large Texts,''
  \textit{arXiv}, 2025.
  \url{https://doi.org/10.48550/arxiv.2505.13208}%
}

\def\bibYeung{%
  R. Yeung and D. Kartsaklis,
  ``A CCG-Based Version of the DisCoCat Framework,''
  \textit{arXiv (Cornell University)}, 2021.
  \url{https://doi.org/10.48550/arxiv.2105.07720}%
}

\def\bibCoeckeSadrzadeh{%
  B. Coecke, M. Sadrzadeh, and S. Clark,
  ``Mathematical Foundations for a Compositional Distributional Model of Meaning,''
  \textit{arXiv (Cornell University)}, vol. 36, no. 1, pp. 345--384, 2010.
  \url{https://doi.org/10.48550/arxiv.1003.4394}%
}

\def\bibMurauer{%
  J. Murauer,
  ``Analyzing Word Predictions by Quantum Natural Language Processing,''
  Master's thesis, Ludwig-Maximilians-Universit\"{a}t M\"{u}nchen, 2023.%
  \url{https://bib.nm.ifi.lmu.de/pdf/mura23.pdf}
}

\def\bibKhatri{%
  N. Khatri,
  ``Experimental Comparison of Ans\"{a}tze for Quantum Natural Language Processing,''
  Master's thesis, University of Oxford, 2022.
  \url{https://www.cs.ox.ac.uk/people/aleks.kissinger/theses/khatri-thesis.pdf}%
}

\def\bibToumi{%
  A. Toumi and G. De Felice,
  ``Higher-Order DisCoCat (Peirce-Lambek-Montague Semantics),''
  \textit{Electronic Proceedings in Theoretical Computer Science}, vol. 429, pp. 130--145, 2025.
  \url{https://doi.org/10.4204/eptcs.429.7}%
}

\def\bibLorenz{%
  R. Lorenz, A. Pearson, K. Meichanetzidis, D. Kartsaklis, and B. Coecke,
  ``QNLP in Practice: Running Compositional Models of Meaning on a Quantum Computer,''
  \textit{Journal of Artificial Intelligence Research}, vol. 76, pp. 1305--1342, 2023.
  \url{https://doi.org/10.1613/jair.1.14329}%
}

\def\bibGuarasci{%
  R. Guarasci, G. De Pietro, and M. Esposito,
  ``Quantum Natural Language Processing: Challenges and Opportunities,''
  \textit{Applied Sciences}, vol. 12, no. 11, p. 5651, 2022.
  \url{https://doi.org/10.3390/app12115651}%
}

\def\bibYi{%
  H. Yi, K. Ban, M. Park, and K. Kong,
  ``Quantum Integration Networks for Efficient Monte Carlo in High-Energy Physics,''
  \textit{arXiv (Cornell University)}, 2025.
  \url{https://doi.org/10.48550/arxiv.2510.10501}%
}

\def\bibWang{%
  S. Wang, P. Czarnik, A. Arrasmith, M. Cerezo, L. Cincio, and P. J. Coles,
  ``Can Error Mitigation Improve Trainability of Noisy Variational Quantum Algorithms?''
  \textit{arXiv (Cornell University)}, 2021.
  \url{https://doi.org/10.48550/arxiv.2109.01051}%
}

\def\bibChu{%
  C. Chu, Q. Lou, F. Chen, and L. Jiang,
  ``QNBAD: Quantum Noise-induced Backdoor Attacks against Zero Noise Extrapolation,''
  \textit{Proceedings of the Network and Distributed System Security (NDSS) Symposium}, 2026.
  \url{https://dx.doi.org/10.14722/ndss.2026.241665}%
}

\def\bibNjiki{%
  J. R. M. Njiki, N. Innan, A. Marchisio, M. Kashif, J. Dricot, and M. Shafique,
  ``Robustness Evaluation of Hybrid Quantum Neural Networks under Noise Models via System-Level Error Mitigation,''
  \textit{arXiv (Cornell University)}, 2026.
  \url{https://doi.org/10.48550/arxiv.2604.17515}%
}

\def\bibRusso{%
  V. Russo,
  ``ZNE Technique Details,''
  \textit{The QEM Zoo}.
  \url{https://qemzoo.com/technique.html?id=zne}%
}

\def\bibSayapin{%
  Z. Sayapin, D. Rabinovich, N. Korolev, and K. Lakhmanskiy,
  ``Zero-Noise Extrapolation via Cyclic Permutations of Quantum Circuit Layouts,''
  \textit{arXiv (Cornell University)}, 2025.
  \url{https://doi.org/10.48550/arxiv.2511.02901}%
}

\def\bibPennyLaneZNE{%
  Xanadu,
  ``catalyst.mitigate\_with\_zne --- Catalyst Documentation,''
  \textit{PennyLane Documentation}, n.d.
  \url{https://docs.pennylane.ai/projects/catalyst/en/stable/code/api/catalyst.mitigate_with_zne.html}%
}

\def\bibKoenig{%
  K. F. Koenig, F. Reinecke, W. Hahn, and T. Wellens,
  ``Inverted-Circuit Zero-Noise Extrapolation for Quantum-Gate Error Mitigation,''
  \textit{Physical Review A}, vol. 110, no. 4, 2024.
  \url{https://doi.org/10.1103/PhysRevA.110.042625}%
}

\def\bibIBMErrorMitigation{%
  IBM Quantum,
  ``Error mitigation and suppression techniques,''
  \textit{IBM Quantum Documentation}.
  \url{https://quantum.cloud.ibm.com/docs/en/guides/error-mitigation-and-suppression-techniques}%
}

\def\bibLiu{%
  J. Liu and X. Wang,
  ``Accelerating Noisy Variational Quantum Algorithms with Physics-Informed Denoising Networks,''
  \textit{arXiv (Cornell University)}, 2026.
  \url{https://doi.org/10.48550/arxiv.2605.02066}%
}

\def\bibMohammadipour{%
  P. Mohammadipour and X. Li,
  ``Direct Analysis of Zero-Noise Extrapolation: Polynomial Methods, Error Bounds, and Simultaneous Physical-Algorithmic Error Mitigation,''
  \textit{Quantum}, vol. 9, p. 1909, 2025.
  \url{https://doi.org/10.22331/q-2025-11-14-1909}%
}

% -------------------------------------------------------
% Document
% -------------------------------------------------------
\begin{document}

\maketitle

\begin{abstract}
  \myabstract
\end{abstract}

\begin{IEEEkeywords}
  \mykeywords
\end{IEEEkeywords}

% -------------------------------------------------------
% Section I: Introduction
% -------------------------------------------------------
\section{Introduction}
\label{sec:introduction}
Research in Natural Language Processing (NLP) increasingly incorporates quantum-enhanced machine learning methods. 
However, current Noisy Intermediate-Scale Quantum (NISQ) devices remain highly sensitive to 
decoherence and gate errors, this noise, limits their reliability. \cite{b_phukan2025}.
% Monte Carlo simulations were conducted to observe how increasing noise and changes in measurement precision impact 
% classification accuracy in a constrained quantum circuit.

NISQ-era quantum computers face two main technical challenges: the rapid environmental decoherence 
and the gate error rates, which depending on the specific hardware platform and experimental setup;
are often reported in the 0.1\% to 1\% ranges. 
Compared to classical computer systems, quantum noise compounds dynamically with every physical operation, this
causes the fidelity of the overall quantum state to decay as circuit depth and quantum gate counts scale. 
This operational decay often restricts reliable computation to shallow architectures of no more than 100 layers \cite{b_zhu2026}.
Without error mitigation strategies, scaling current hardware beyond architectural limits remains constrained.

Optimizing high-dimensional Natural Language Understanding (NLU) models often
introduces the additional challenge of the 'barren plateau' phenomenon. 
The barren plateau phenomenon is where cost function gradients vanish exponentially as 
the qubit count scales; this amplifies the hardware instability \cite{b_phukan2025}. 
These plateaus effectively trap optimizers in suboptimal search regions, often leading to premature 
convergence and poor classification performance of Machine Learning Systems \cite{b_lung2026}

Previous analyses of quantum noise dynamics have shown that controlled stochastic 
effects can behave like a regularizer, leading to reductions in both the model’s Lipschitz 
constant and the Quantum Fisher Information Matrix (QFIM). \cite{b_zhu2026}.
These mechanisms bring flatter minima in the loss landscape, which have been shown to yield tighter
Probably Approximately Correct (PAC)-Bayesian generalization bounds \cite{b_xhelollari}.

However, this benefit introduces a critical tension between regularizing the model and maintaining its trainability.\cite{b_zhu2026, b_cincio2021}.
If pushed beyond a critical threshold, this state-space contraction triggers noise-induced barren plateaus (NIBPs), 
where gradients vanish exponentially and beneficial stochastic fluctuations become disruptive instead of
aiding optimization \cite{b_xhelollari, b_higgs}.How the model responds to the noise depends 
strictly on initialization: it assists under-optimized models 
but disrupts well-converged ones, establishing a negative correlation with baseline performance. \cite{b_zhu2026}.

A Monte Carlo simulation framework isolates the transition between beneficial regularization 
and catastrophic decoherence. The pipeline evaluates 1,000 independent trials for each discrete 
configuration of depolarizing noise (1–10\%) and measurement shot counts. 
The high-volume sampling extracts the statistical variance introduced by hardware fluctuations 
and finite-shot sampling noise, both of which actively degrade gradient estimation and model output.


% A system-wide reliability boundary provides a framework to evaluate where hardware-induced 
% decoherence begins to overwhelm predictive utility. This shift typically manifests as noise 
% introduces granularity and local irregularities into the continuous decision function. \cite{b_seilkhan2026}

% -------------------------------------------------------
% Section II: Related Literature
% -------------------------------------------------------
\section{Related Literature}
\label{sec:related_literature}

\subsection{The Convergence of Quantum Computing and NLP}

The convergence of quantum computing and Natural Language Processing (NLP) represents a fundamental paradigm shift 
from purely statistical, associative methods to structure-preserving, 
categorical representations of meaning \cite{b_phukan2025, b_pal2025}.
This section focuses on the theoretical foundations of this convergence, 
specifically the "Quantum Native" hypothesis and the principle of compositionality.

\subsubsection{The "Quantum Native" Hypothesis}

The foundational motivation for Quantum Natural Language Processing (QNLP) is the proposition that 
human language is "quantum native" \cite{b_phukan2025, b_pal2025, b_coecke2020}.
The 'Quantum Native' hypothesis suggests that natural language fits naturally into quantum 
computers because the mathematical structures of grammar match those of quantum mechanics. 
Both frameworks use vector spaces to represent individual elements and tensor products to 
calculate how these elements interact.

This shared mathematical framework allows for a direct and systematic translation of a sentence’s 
grammatical string diagram into a quantum circuit, where words act as quantum states and 
grammatical dependencies are represented as entangling operations \cite{b_pal2025, b_coecke2020}.
Because language and quantum mechanics share this underlying categorical logic, quantum computers
are viewed as the "native hardware" for executing linguistic models.

A key advantage of this approach in the current Noisy Intermediate-Scale Quantum (NISQ) era is 
that linguistic structure is essentially provided as a "free lunch," meaning that the grammatical rules of 
a sentence are encoded as the hardware's physical configuration rather than being learned from data \cite{b_coecke2020}.
While classical systems view the encoding of complex grammar as exponentially expensive, the variational 
quantum circuit paradigm allows syntactic structure to be provided for "free" by defining the quantum circuit skeleton 
before training begins \cite{b_phukan2025, b_coecke2020}.

In classical models, such as Large Language Models (LLMs), grammatical rules must be learned from scratch through 
massive parameter budgets, often resulting in "black box" systems whose internal reasoning is opaque and lacks structural 
interpretability \cite{b_phukan2025, b_krawchuk2025}.

To expound on this empirical evidence, the QSANN 
(Quantum Self-Attention Neural Network) model serves as a primary 
example of this efficiency, particularly in text classification tasks.
On sentiment classification benchmarks using real-world data like the Yelp dataset, 
the QSANN was shown to outperform comparable classical self-attention baselines while using only 
49 quantum parameters against the 785 classical parameters required by the classical model—a reduction of >90\% \cite{b_phukan2025}.

This drastic reduction is rooted in the exponential representational capacity of quantum systems. 
Theoretically, an $n$-qubit system can represent $2^{n}$ states in parallel, allowing quantum 
models to inhabit high-dimensional Hilbert spaces that are exponentially larger 
than the latent spaces accessible to classical models of equivalent size. \cite{b_zengcoecke2016, b_pal2025}

\subsection{Compositionality and the DisCoCat Framework}
The theoretical cornerstone of modern Quantum Natural Language 
Processing (QNLP) is the Distributional Compositional Categorical (DisCoCat) framework \cite{b_yeung2021}. 
This model provides a mathematically rigorous unification of two central pillars of
linguistic theory: the distributional hypothesis, which posits that a word’s meaning is
determined by its context, and the principle of compositionality, which states that the meaning of a whole 
expression is a function of the meanings of its parts and the rules used to combine them \cite{b_coeckesadrzadeh2010}.

\subsubsection{Mathematical Mechanics of DisCoCat}
The DisCoCat (Distributional Compositional Categorical) framework operates by 
establishing a formal, structure-preserving mapping between formal grammar and vector space semantics \cite{b_yeung2021, b_murauer2023}.
This unification allows the model to leverage the quantitative strengths of the 
distributional hypothesis — where word meaning is derived from context, and the qualitative rigor 
of the principle of compositionality, which dictates how those meanings combine according to grammatical rules \cite{b_murauer2023, b_khatri2022}.

\subsubsection{Pregroup Grammar}
On the syntactic side, the DisCoCat framework primarily utilizes pregroup grammars, 
an algebraic approach introduced by Joachim Lambek where words are assigned abstract types.
A pregroup is a partially ordered monoid where every element $p$ has a left adjoint $p^l$
and a right adjoint $p^r$. These adjoints act as directional inverses for grammatical dependencies \cite{b_yeung2021, b_coeckesadrzadeh2010, b_khatri2022}.
For instance:

\begin{itemize}
  \item A \textbf{noun} is assigned a basic type $n$.
  \item A \textbf{transitive verb} is assigned a compound type $n^r s n^l$, indicating it requires a subject ($n$) on its left and an object ($n$) on its right to form a well-typed sentence of type $s$.
  \item Linguistic reductions, such as $n \cdot n^r \rightarrow 1$, are represented as ``cups'' in a string diagram, signaling the fulfillment of these grammatical requirements.
\end{itemize}

\subsubsection{Hilbert Space Representations}
Under this framework, word meanings are represented as vectors or higher-order tensors within high-dimensional Hilbert spaces.
Words with basic types are represented as vectors, while functional words like verbs are represented
as tensors whose order is determined by their compound types \cite{b_yeung2021, b_murauer2023, b_toumi2025}. Consequently, the process of parsing a sentence 
becomes a physical process of composing these tensors \cite{b_lorenz2023, b_guarasci2022}.
This "quantum-native" structure allows the grammatical derivation of a sentence to
dictate the exact wiring of a quantum circuit, where word meanings are instantiated as parameterized quantum operations.

\subsection{Dynamics and Duality of Quantum Noise in NISQ Systems}
The utility of Noisy Intermediate-Scale Quantum (NISQ) devices is fundamentally governed by
the interplay between hardware-induced disturbances and the model's Hilbert Space. Unlike
classical noise, quantum noise is categorized by complex decoherence mechanisms primarily
\textbf{energy relaxation ($T_1$)} and \textbf{dephasing ($T_2$)} which manifest at the
algorithmic level as depolarizing, amplitude damping, or bit-flip error channels \cite{b_zhu2026, b_yi2025}.

\subsubsection{Principles and Mechanics of Zero-Noise Extrapolation (ZNE)}
To mitigate the systematic biases introduced by hardware decoherence, Zero-Noise Extrapolation (ZNE) 
has emerged as a premier hardware-agnostic post-processing strategy for the NISQ era \cite{b_wang2021, b_chu2026, b_njiki2026}.
ZNE operates on the core principle that if the expectation value of a physical 
observable varies smoothly with the system's noise rate, one can estimate the ideal 
"zero-noise" result by evaluating the circuit at multiple elevated noise levels and extrapolating 
the results back to the noise-free limit \cite{b_russo, b_sayapin2025, b_pennylane_zne}.
The workflow generally consists of two primary stages: noise amplification and extrapolation.

\paragraph{Controlled Noise Amplification}

In the first stage, the effective noise strength ($\lambda$) is increased in a 
controlled manner without altering the logical state of the computation. 
Several methodologies exist for scaling noise:


\begin{enumerate}
  \item \textbf{Digital Unitary Folding}: This gate-level technique replaces a unitary $U$ with the logically equivalent 
  sequence $U(U^\dagger U)^n$. This increases the circuit depth and time-to-solution, 
  effectively scaling the cumulative noise by odd factors of $s = 2n + 1$. \cite{b_sayapin2025, b_koenig2024}.
  \item \textbf{Analog Pulse Stretching:} This approach involves physically
  stretching the duration of microwave control pulses while
  inversely scaling their amplitude, exposing qubits to
  environmental decoherence mechanisms ($T_1$ and $T_2$) for
  longer durations. \cite{b_russo, b_koenig2024, b_liu2026}
  \item \textbf{Probabilistic Error Amplification (PEA)}: A model-based approach where a twirled noise model is learned
  for each circuit layer and then amplified by inserting random Pauli gates\cite{b_koenig2024}.
  \item \textbf{Cyclic Layout Permutations (CLP-ZNE)}: This "gray-box" method exploits the spatial 
  non-uniformity of gate fidelities on a chip. By executing the logical circuit across different physical qubit mappings 
  (layouts), it effectively scales the noise without adding extra gates or depth. \cite{b_sayapin2025}
\end{enumerate}


\paragraph{Extrapolation to the Zero-Noise Limit}
Once expectation values $C_\lambda(\boldsymbol{\theta})$ are collected at multiple noise factors
(typically $\lambda = 1, 3, 5$), a classical fitting model is applied. Common
fitting functions include linear, polynomial, exponential, or Richardson
extrapolation. The mitigated expectation value is determined by
taking the limit as the noise scale factor approaches zero: $C_{\text{ZNE}}(\boldsymbol{\theta}) =
\lim_{\lambda \to 0} C_\lambda(\boldsymbol{\theta})$ \cite{b_chu2026, b_njiki2026, b_russo, b_liu2026, b_mohammadipour2025}.

\subsection{Synthesis and Research Gap}

The reviewed studies show that Quantum Natural Language Processing (QNLP) can represent linguistic
structure through quantum circuits and compositional models \cite{b_phukan2025, b_coecke2020, b_lorenz2023}.
Related works also show that noise affects variational quantum models in two ways: it can sometimes
support generalization, but it can also reduce trainability when the noise becomes too strong
\cite{b_zhu2026, b_xhelollari, b_cincio2021}. Meanwhile, Zero-Noise Extrapolation (ZNE) is commonly
used as an error mitigation method for estimating circuit outputs closer to the zero-noise case
\cite{b_wang2021, b_russo, b_sayapin2025, b_mohammadipour2025}.

Based on these studies, this paper focuses on a smaller and more controlled setting: a two-qubit
Variational Quantum Classifier for binary intent classification. The gap addressed in this study is
the limited evaluation of how depolarizing noise, measurement shot counts, and ZNE affect the accuracy
and variance of a simple VQC model under repeated Monte Carlo trials.

\begin{thebibliography}{00}
  \bibitem{b_phukan2025}\bibPhukan %A survey of Quantum Natural Language Processing: From Compositional models to NISQ-Era Empiricism
  \bibitem{b_zhu2026}\bibZhu %Rethinking Quantum Noise in Quantum Machine Learning: When Noise Improves Learning
  \bibitem{b_lung2026}\bibLung %A Noisy Optimization Mechanism for Variational Quantum Classifiers
  \bibitem{b_xhelollari}\bibXhelollari %Robust Hybrid Quantum-Classical Latent Diffusion Models via Quantum Noise
  \bibitem{b_cincio2021}\bibCincio %Machine Learning of Noise-Resilient Quantum Circuits
  \bibitem{b_higgs}\bibHiggs %Impact of Circuit Depth versus Qubit Count on Variational Quantum Classifiers for Higgs Boson Signal Detection
  \bibitem{b_seilkhan2026}\bibSeilkhan %Comparing Classical and Quantum Variational Classifiers on the XOR Problem
  \bibitem{b_zengcoecke2016}\bibZengCoecke %Quantum Algorithms for Compositional Natural Language Processing
  \bibitem{b_pal2025}\bibPal %Proceedings of the QuantumNLP: Integrating Quantum Computing with Natural Language Processing
  \bibitem{b_coecke2020}\bibCoecke %Foundations for Near-Term Quantum Natural Language Processing
  \bibitem{b_krawchuk2025}\bibKrawchuk %Efficient Generation of Parameterised Quantum Circuits from Large Texts
  \bibitem{b_yeung2021}\bibYeung %A CCG-Based Version of the DisCoCat Framework
  \bibitem{b_coeckesadrzadeh2010}\bibCoeckeSadrzadeh %Mathematical Foundations for a Compositional Distributional Model of Meaning
  \bibitem{b_murauer2023}\bibMurauer %Analyzing Word Predictions by Quantum Natural Language Processing Master's Thesis
  \bibitem{b_khatri2022}\bibKhatri %Experimental Comparison of Ansätze for Quantum Natural Language Processing
  \bibitem{b_toumi2025}\bibToumi %Higher-Order DisCoCat (Peirce-Lambek-Montague Semantics)
  \bibitem{b_lorenz2023}\bibLorenz %QNLP in Practice: Running Compositional Models of Meaning on a Quantum Computer
  \bibitem{b_guarasci2022}\bibGuarasci %Quantum Natural Language Processing: Challenges and Opportunities
  \bibitem{b_yi2025}\bibYi %Quantum Integration Networks for Efficient Monte Carlo in High-Energy Physics
  \bibitem{b_wang2021}\bibWang %Can Error Mitigation Improve Trainability of Noisy Variational Quantum Algorithms?
  \bibitem{b_chu2026}\bibChu %QNBAD: Quantum Noise-induced Backdoor Attacks against Zero Noise Extrapolation
  \bibitem{b_njiki2026}\bibNjiki %Robustness Evaluation of Hybrid Quantum Neural Networks under Noise Models via System-Level Error Mitigation
  \bibitem{b_russo}\bibRusso %ZNE Technique Details - The QEM Zoo
  \bibitem{b_sayapin2025}\bibSayapin %Zero-Noise Extrapolation via Cyclic Permutations of Quantum Circuit Layouts
  \bibitem{b_pennylane_zne}\bibPennyLaneZNE %catalyst.mitigate_with_zne documentation
  \bibitem{b_koenig2024}\bibKoenig %Inverted-circuit zero-noise extrapolation for quantum-gate error mitigation
  \bibitem{b_ibm_error_mitigation}\bibIBMErrorMitigation %Error mitigation and suppression techniques
  \bibitem{b_liu2026}\bibLiu %Accelerating Noisy Variational Quantum Algorithms with Physics-Informed Denoising Networks
  \bibitem{b_mohammadipour2025}\bibMohammadipour %Direct analysis of Zero-Noise extrapolation: polynomial methods, error bounds, and simultaneous Physical-Algorithmic error mitigation
\end{thebibliography}

\end{document}